\section{The dynamic layer: finite review envelopes}\label{sec:dynamic}

The static theory certifies that an atom has no common safe action \emph{now}. In a dynamic problem the controller acts, observes, and acts again; the atom is replaced by an information state, the action by a policy tree, and the obstruction by a value. This section develops the finite case completely and exactly: finite scenario spaces, finitely many review stages, finite action sets and observation cells that may depend on earlier actions. Everything is elementary, and that is the point---the identities that are easy to get wrong in continuous formulations (which conditioning is legitimate, when a prefix may be committed, what the adversary's choice of a coupling does) become finite checks.

\subsection{Review trees and the envelope}

\begin{definition}[finite review problem]\label{def:tree}
A finite review problem consists of a finite scenario space $\Omega$ with a law $\mu$ of full support; $K$ stages with finite action sets $A_1,\dots,A_K$; observation maps $o_k:\Omega\times A_1\times\dots\times A_{k-1}\to O_k$, where $o_k(\omega,a_{<k})$ is what the controller sees before choosing $a_k$ and may depend on its earlier actions; and a terminal functional $G:\Omega\times A_1\times\dots\times A_K\to\Qr$. A \emph{history} at stage $k$ is $h_k=(o_1,a_1,\dots,a_{k-1},o_k)$ and its \emph{cell} $C(h_k)\subset\Omega$ is the set of scenarios producing it. An \emph{adapted policy tree} $\pi$ assigns an action $\pi(h_k)\in A_k$ to every reachable history; it generates the action path $\pi(\omega)$ and the value $J(\pi)=\E_\mu[G(\omega,\pi(\omega))]$. The \emph{node values} are defined backwards: for a reachable history $h_k$ and $a\in A_k$,
\begin{equation}\label{eq:N}
Q_k(h_k,a)=\begin{cases}\E_\mu\bigl[G(\omega,a_{<k},a)\,\big|\,C(h_k)\bigr], & k=K,\\[2pt]
\displaystyle\sum_{o\in O_{k+1}}\mu\bigl(C(h_k,a,o)\,\big|\,C(h_k)\bigr)\,N_{k+1}(h_k,a,o), & k<K,\end{cases}\qquad
N_k(h_k)=\max_{a\in A_k}Q_k(h_k,a),
\end{equation}
and the \emph{envelope} is $Q_0=\sum_{o\in O_1}\mu(C(o))\,N_1(o)$.
\end{definition}

\begin{theorem}[the envelope is the tree supremum]\label{thm:envelope}
For every finite review problem: (i) $Q_0=\max_\pi J(\pi)$ over adapted policy trees, and any tree that selects a maximiser of $Q_k(h_k,\cdot)$ at every reachable history attains it; (ii) for every fixed tree $\pi$ the process $M_k=\E_\mu[G(\omega,\pi(\omega))\mid h_k^\pi]$ along the histories generated by $\pi$ satisfies $M_k=\E_\mu[M_{k+1}\mid h_k^\pi]$ (martingale property) and $J(\pi)=\E_\mu[M_1]$ (tower identity); (iii) refining an observation map (replacing $o_k$ by a map whose cells are subsets of the original cells) cannot decrease $Q_0$.
\end{theorem}

\begin{proof}
(ii) is the tower property of conditional expectation on the finite filtration generated by $\pi$'s histories; the cells of $h^\pi_{k+1}$ partition the cell of $h^\pi_k$ because the next action is determined by $h_k^\pi$. (i) By backward induction on $k$: for every tree $\pi$ and reachable $h_k$, $\E_\mu[G^\pi\mid C(h_k)]\le N_k(h_k)$, with equality when $\pi$ selects maximisers from $h_k$ onwards; the induction step is $\E[G^\pi\mid C(h_k)]=\sum_o\mu(C(h_k,\pi(h_k),o)\mid C(h_k))\,\E[G^\pi\mid C(h_k,\pi(h_k),o)]\le Q_k(h_k,\pi(h_k))\le N_k(h_k)$. Summing over the initial cells gives $J(\pi)\le Q_0$ with equality for the greedy tree. (iii) Every tree adapted to the coarser map is adapted to the finer one, so the supremum in (i) can only grow.
\end{proof}

The theorem is standard dynamic programming for a finite partially observed problem \cite{Astrom1965,BertsekasShreve1978}; its value here is that it is stated on trees, where three consequences become exact statements that are verified by enumeration on every instance of the paper (Section~\ref{sec:verification}).

\begin{corollary}[committed prefixes]\label{cor:prefix}
For every reachable history $h_k$ and action $a$, the supremum of $\E_\mu[G^\pi\mid C(h_k)]$ over trees that play $a$ at $h_k$ equals $Q_k(h_k,a)$. In particular the value of a review scheme under a first action committed in advance is the node value of that action, and the loss of commitment is $N_1-Q_1(\cdot,a)$.
\end{corollary}

\begin{example}[posterior-only evaluation is not an envelope]\label{ex:posterior}
A blind first action $a_1\in\{0,1\}$ is followed by the revelation of a fair bit $s$; the reward is $\ind{a_1=s}$. The adapted value is $Q_0=\tfr12$. Evaluating the first action under the \emph{terminal} posterior---$\E[\max_{a_1}\E[\ind{a_1=s}\mid s]]=1$---conditions an earlier action on a later observation and overstates the value by $\tfr12$. The recursion \eqref{eq:N} cannot make this error because $Q_1(h_1,a)$ averages over the cells of the \emph{next} observation given the current one. Any formulation that reasons with the posterior at the end of the horizon alone (as a ``residue'' of the information) must be checked against \eqref{eq:N}; on this instance the discrepancy is maximal.
\end{example}

\subsection{The two-window instance}\label{sec:twowindow}

A hidden sign $s\in\{\pm1\}$ and two shocks $d_1,d_2\in\{\pm1\}$ are uniform and independent under the product prior $\mu$. A scalar state starts at $z_0\in\Qr$ and moves as $z_1=z_0+su_1+d_1$, $z_2=z_1+su_2+d_2$ with actions $u_k\in\{-2,0,2\}$; the first action is blind, the second sees $z_1$; the terminal functional is the corridor margin $G=2-|z_2|$. Because the sign of the action's effect is unknown, the first action is a probe as much as a control: with $u_1=\pm2$ the four values of $z_1$ are distinct and reveal $s$; with $u_1=0$ nothing is learned about $s$. The reachable histories carry $171=81+9+81$ adapted trees.

At $z_0=2$ the envelope is $Q_0=\tfr12$ with first-stage values $Q_1(\cdot,-2)=\tfr12$, $Q_1(\cdot,0)=0$, $Q_1(\cdot,+2)=\tfr12$; after $u_1=+2$ the review cells $z_1\in\{-1,1,3,5\}$ have node values $(1,1,1,-1)$, with ties $u_2\in\{0,-2\}$ at $z_1=-1$ and $\{0,+2\}$ at $z_1=1$ and the unique brake $u_2=-2$ at $z_1\in\{3,5\}$ (Figure~\ref{fig:review}a). Eight of the $171$ trees attain $\tfr12$; the tower and martingale identities hold on all $171$; the open-loop supremum is $0$, attained by $(-2,+2)$, $(0,0)$ and $(+2,-2)$, so the review is worth exactly $\tfr12$.

\begin{figure}[t]
\centering
\includegraphics[width=\textwidth]{fig3_review.pdf}
\caption{The two-window instance (Theorem~\ref{thm:envelope}, Proposition~\ref{prop:robust}, Example~\ref{ex:sweep}). (a) The review tree at $z_0=2$: first-stage values, and the four review cells after $u_1=+2$ with their node values and optimal second actions. (b) Envelope $Q_0$, open-loop supremum and the robust values $R^\pm$ over the coupling class $\calM$ as functions of the initial distance $z_0$; the shaded region $z_0\ge4$ is the dynamically certified obstruction, where every adapted tree is defeated with positive probability.}
\label{fig:review}
\end{figure}

\subsection{Robust values over a class of couplings}

The product prior is one coupling of the one-step marginals; a robust reading lets an adversary choose the coupling. Let $\calM\subset\calP(\Omega)$ be a polytope of laws containing $\mu$ (in the instance: all laws on $\{\pm1\}^3$ with uniform coordinate marginals, a polytope of dimension $4$ with six vertices, whose barycentre is the product prior).

\begin{proposition}[robust envelope: weak duality and its closure]\label{prop:robust}
Let $\Pi$ be the finite set of adapted trees and define the lower and upper robust values
\begin{equation}\label{eq:robust}
R^-=\max_{\pi\in\Pi}\ \min_{\mu'\in\calM}\E_{\mu'}[G^\pi],\qquad R^+=\min_{\mu'\in\calM}\ \max_{\pi\in\Pi}\E_{\mu'}[G^\pi].
\end{equation}
Then (i) $R^-\le R^+\le Q_0(\mu)$ for every $\mu\in\calM$; (ii) $R^+=\max_{w\in\Delta(\Pi)}\min_{\mu'\in\calM}\sum_\pi w_\pi\E_{\mu'}[G^\pi]$: mixing over trees closes the gap, and a pair $(\mu^\ast,w^\ast)$ with $\max_\pi\E_{\mu^\ast}G^\pi=\min_{\mu'\in\calM}\sum_\pi w^\ast_\pi\E_{\mu'}G^\pi$ certifies $R^+$ exactly; (iii) $R^+<0$ iff some coupling $\mu^\ast\in\calM$ defeats every tree in expectation, in which case under $\mu^\ast$ every tree is defeated on a set of positive $\mu^\ast$-measure; the pair ($\mu^\ast$, the recursion \eqref{eq:N} under $\mu^\ast$) is a dynamic certificate of obstruction checkable by one backward induction; (iv) all four quantities $R^\pm$, $Q_0$, and the certifying pair are rational when the data are, being optimal values and vertices of linear programmes.
\end{proposition}

\begin{proof}
(i) is weak duality (for every $\pi$ and $\mu'$, $\min_{\mu''}\E_{\mu''}G^\pi\le\E_{\mu'}G^\pi\le\max_{\pi'}\E_{\mu'}G^{\pi'}$) and $\mu\in\calM$ with Theorem~\ref{thm:envelope}(i). (ii) The payoff $(\mu',w)\mapsto\sum_\pi w_\pi\E_{\mu'}G^\pi$ is bilinear on the product of two polytopes, so von Neumann's theorem \cite{vonNeumann1928} gives $\min_{\mu'}\max_w=\max_w\min_{\mu'}$, and $\max_w$ of a linear function over the simplex is attained at a vertex, i.e. $\min_{\mu'}\max_w=R^+$. (iii) The first statement is the definition; if $\E_{\mu^\ast}G^\pi<0$ then $\mu^\ast\{G^\pi<0\}>0$. (iv) The inner minima are linear programmes over $\calM$ and the outer problems are linear programmes in $(\mu',t)$ or $(w,t)$.
\end{proof}

\begin{example}[initial distance, couplings and a certified dynamic obstruction]\label{ex:sweep}
Table~\ref{tab:sweep} lists, for the two-window instance and $z_0=0,\dots,8$, the envelope $Q_0$ under the product prior, the open-loop supremum, and $R^\pm$ over the class $\calM$ of couplings with uniform coordinate marginals; every entry is exact and $R^+$ is certified by a pair $(\mu^\ast,w^\ast)$ as in Proposition~\ref{prop:robust}(ii).
\begin{enumerate}[leftmargin=1.6em,itemsep=1pt]
\item \emph{Value of the review.} $Q_0-\text{open-loop}=0,\tfr14,\tfr12,1,\tfr32,\tfr74,2,2,2$: the review is worth nothing at $z_0=0$ and two units of margin from $z_0=6$ on.
\item \emph{The worst coupling can be interior.} At $z_0=2$, $R^+=\tfr12=Q_0$ and the worst coupling is the product prior itself, while each of the six extreme couplings of $\calM$ gives the controller strictly more ($2,2,1,2,1,2$). At $z_0=1$ and $z_0=5$ a non-product coupling is strictly worse than the prior ($R^+=\tfr12<\tfr34$ and $R^+=-\tfr32<-\tfr54$); at $z_0=3$ the worst coupling is supported on four of the eight paths and ties with the prior.
\item \emph{Mixing closes the gap.} At $z_0=2$, $R^-=0<R^+=\tfr12$: no single tree guarantees more than $0$ against all couplings, but the uniform mixture of four trees guarantees $\tfr12$. From $z_0=6$ on, $R^-=R^+=Q_0$: one tree is optimal against every coupling.
\item \emph{Dynamic obstruction.} At $z_0=3$ the envelope is $0$ and every tree is defeated with positive probability although some have zero expected loss. At $z_0=4$ every one of the $171$ trees has negative expectation under the prior and is defeated with probability at least $\tfr38$; $R^+=-\tfr12=Q_0$, so the prior is itself a worst coupling and the pair (prior, recursion) is the dynamic certificate of Proposition~\ref{prop:robust}(iii). This is the review-tree counterpart of the static certificate: an explicit measure under which no adapted policy survives.
\item \emph{Coupling selection is real.} With identical one-step marginals and $G=d_1d_2$, the product coupling gives $0$, the comonotone coupling $1$ and their $\rho$-mixture $\rho$: the one-step marginals under-determine the envelope, which is why $\calM$ must be declared.
\end{enumerate}
\end{example}

\begin{table}[t]
\centering\small
\caption{Two-window instance: envelope, open-loop supremum and robust values against the initial distance (Example~\ref{ex:sweep}); all entries exact, $R^+$ certified on both sides.}
\label{tab:sweep}
\begin{tabular}{lccccccccc}
\toprule
$z_0$ & 0 & 1 & 2 & 3 & 4 & 5 & 6 & 7 & 8\\
\midrule
envelope $Q_0$ (product prior) & $1$ & $\tfr34$ & $\tfr12$ & $0$ & $-\tfr12$ & $-\tfr54$ & $-2$ & $-3$ & $-4$\\
open-loop supremum & $1$ & $\tfr12$ & $0$ & $-1$ & $-2$ & $-3$ & $-4$ & $-5$ & $-6$\\
$R^+=\min_{\mu'}\max_\pi$ & $1$ & $\tfr12$ & $\tfr12$ & $0$ & $-\tfr12$ & $-\tfr32$ & $-2$ & $-3$ & $-4$\\
$R^-=\max_\pi\min_{\mu'}$ & $0$ & $0$ & $0$ & $-1$ & $-1$ & $-2$ & $-2$ & $-3$ & $-4$\\
\bottomrule
\end{tabular}
\end{table}

\subsection{The integrated certificate}\label{sec:integrated}

The static certificate pools rows at one instant. For a finite-horizon trajectory problem with a finite set of branches (hidden modes) and an information schedule, the same pooling applies to the terminal constraint slack integrated along the trajectory, and yields a certificate against every policy adapted to the schedule.

\begin{proposition}[integrated certificate with adapted controls]\label{prop:integrated}
Let $L$ be a finite set of branches, $U$ a compact convex control set, $T>0$, and for each $\ell\in L$ let the terminal slack of a constraint be an affine functional of the control path,
\begin{equation}\label{eq:slack}
F_\ell(u_\ell)=\beta_\ell-\int_0^T k_\ell(\sigma)^\top u_\ell(\sigma)\dd\sigma,
\end{equation}
with integrable kernels $k_\ell$. Let $(\calP_\sigma)_{\sigma\in[0,T]}$ be a refining information schedule (partitions of $L$), and call a family $(u_\ell)_{\ell\in L}$ of measurable $U$-valued controls \emph{adapted} if $u_\ell(\sigma)=u_{\ell'}(\sigma)$ whenever $\ell,\ell'$ lie in the same cell of $\calP_\sigma$. Then for every $\lambda\in\Delta(L)$ and every adapted family,
\begin{equation}\label{eq:gammah}
\min_{\ell\in L}F_\ell(u_\ell)\ \le\ \Gamma_h(\lambda)\ :=\ \sum_\ell\lambda_\ell\beta_\ell+\int_0^T\ \sum_{C\in\calP_\sigma}h_U\Bigl(-\sum_{\ell\in C}\lambda_\ell k_\ell(\sigma)\Bigr)\dd\sigma .
\end{equation}
Hence $\Gamma_h(\lambda)<0$ certifies that every adapted policy violates some branch's constraint at $T$. On the three-branch instance of the companion paper \cite{Abaee2026cert}---a planar double integrator $\ddot p=u$ with $U$ the hexagon $\conv\{\pm n_\ell\}$, $n_1=(1,0)$, $n_{2,3}=(-\tfr35,\pm\tfr45)$, initial states $p_\ell(0)=\tfr{34}{25}n_\ell$, $\dot p_\ell(0)=n_\ell$, facets $n_\ell^\top p\le2$, the branch revealed at time $\tau$ and horizon $T=\tau+1$---the kernels are $k_\ell(\sigma)=(T-\sigma)n_\ell$, the margins $\beta_\ell=\tfr{16}{25}-T$, and $\lambda=(\tfr38,\tfr5{16},\tfr5{16})$ gives
\begin{equation}\label{eq:gamma3}
\Gamma_h(\tau)=\tfrac{7}{50}-\tau ,
\end{equation}
so the review is obstructed for every $\tau>\tfr7{50}$; the threshold is sharp by the policy of the companion paper, and $\Gamma_h(\tau)=-\Gamma(\tau)$ in the companion's sign convention, where $\Gamma(\tau)=\tau-\tfr7{50}$ is the value of a minimisation programme.
\end{proposition}

\begin{proof}
$\min_\ell F_\ell\le\sum_\ell\lambda_\ell F_\ell=\sum_\ell\lambda_\ell\beta_\ell-\int_0^T\sum_{C\in\calP_\sigma}\Bigl(\sum_{\ell\in C}\lambda_\ell k_\ell(\sigma)\Bigr)^{\!\top}u_C(\sigma)\dd\sigma$, where $u_C(\sigma)$ is the common control of the cell $C$; and $-g^\top u\le h_U(-g)$ for $u\in U$. For the instance, $\sum_\ell\lambda_\ell n_\ell=0$ kills the integrand on the blind window $[0,\tau)$, where the schedule has the single cell $L$; after $\tau$ the cells are singletons and $h_U(-n_\ell)=1$ because $-n_\ell$ is a vertex of the hexagon and $\|n_\ell\|=1$; $\int_\tau^T(T-\sigma)\dd\sigma=\tfr12$; and $\tfr{16}{25}-(\tau+1)+\tfr12=\tfr7{50}-\tau$. The slack \eqref{eq:slack} is $2-n_\ell^\top p_\ell(T)$ with $n_\ell^\top p_\ell(T)=\tfr{34}{25}+T+\int_0^T(T-\sigma)n_\ell^\top u_\ell(\sigma)\dd\sigma$.
\end{proof}

Two features of \eqref{eq:gammah} matter beyond the instance. First, the support function is evaluated at the \emph{negative} pooled kernel; on the centrally symmetric hexagon the sign is immaterial, but on a translated control set the two conventions differ (check G1.3), and the certificate of an asymmetric actuator must use \eqref{eq:gammah} as written. Second, the identity reproduces the one-step theory: a schedule with a single instant and a single cell is Theorem~\ref{thm:dual}'s weak direction with $k_\ell$ the drift coefficients. The converse---that $\Gamma_h<0$ for some $\lambda$ whenever every adapted policy fails---is the finite-horizon bridge theorem of the companion paper \cite{Abaee2026cert}, proved there by a mesh-exact linear programme; it is not re-proved here.

\begin{figure}[t]
\centering
\includegraphics[width=\textwidth]{fig4_helly_gamma.pdf}
\caption{(a) Minimal infeasible subfamily sizes of random row systems with $m=k+3$ rows on $[-1,1]^k$ (Remark~\ref{rem:random}; $n$ is the number of infeasible systems among the draws). The size $k+1$ of Theorem~\ref{thm:helly} is attained but rare; with rank-one coefficients in $k=3$ the size never exceeds $2$. (b) The integrated certificate of the three-branch instance, $\Gamma_h(\tau)=\tfr7{50}-\tau$ (Proposition~\ref{prop:integrated}), against the companion's $\Gamma(\tau)=\tau-\tfr7{50}$; both cross zero at $\tau^\ast=\tfr7{50}$.}
\label{fig:helly}
\end{figure}
